
% Authoritative LaTeX manuscript; Markdown is a generated compatibility projection.

\emph{What should the environment learn from this recurrence?}

When the same kind of failure recurs, ask whether the environment should change. First classify what keeps going wrong: missing knowledge, a weak representation, an implicit obligation, late evidence, misplaced enforcement, a failed mechanism, an architectural weakness, or an obsolete control. Then decide whether the lesson is worth making durable.

Engineering Capital asks whether existing structure still earns its keep; Governance Conversion asks whether a recurring lesson should change that structure.

The conversion decision has three stages:

\begin{itemize}
\tightlist
\item
  \textbf{Classify the gap.} Identify the recurring problem rather than its latest symptom: missing knowledge, representation, obligation, evidence, enforcement, mechanism, architecture, or an obsolete control.
\item
  \textbf{Choose the durable response.} Possible responses include a representation, constraint, sensor, validator, gate, procedure, architectural change, retirement of stale machinery, or deliberate retention of human judgment.
\item
  \textbf{Price the investment.} Build the response only when expected recurrence cost and consequence justify both construction and carrying cost.
\end{itemize}

% mage-code-language: text
\begin{verbatim}
GOVERNANCE CONVERSION
"What should the environment learn?"
Recurring event
      ↓
CLASSIFY THE GAP
knowledge · representation · obligation · evidence
enforcement · mechanism · architecture · obsolete control
      ↓
CHOOSE THE RESPONSE
model · constraint · sensor · validator · gate
procedure · redesign · retire · keep human
      ↓
PRICE THE INVESTMENT
future value > construction + carrying cost ?
\end{verbatim}

The goal is fewer costly recurrences and durable responses that repay their carrying cost. Cheap, rare failures may remain unconverted.

See \href{3.4-governance-conversion.html}{\magevalue{sec:governance-conversion}} for the full treatment.
